% !TeX root = ../main.tex
\section{引言}
\label{sec:introduction}
随着 AI 模型规模与计算需求的增长，低精度浮点计算正成为提升计算效率的重要途径，FP8 及更低位宽格式已逐步用于实际模型。DeepSeek-V3 在混合精度训练中以 FP8 执行主要线性层的前向与反向矩阵乘法~\cite{deepseekv3}；GLM-4.5 在保留 BF16 参数更新的同时，通过在线 FP8 量化加速强化学习中的采样推理~\cite{glm45}。面向更低位宽，DeepSeek-V4 通过量化感知训练将 MoE 专家权重量化为 MXFP4 用于部署，并在长上下文注意力索引器中以 MXFP4 存储 Q/K 激活、执行点积计算~\cite{deepseekv4}。这些实践表明，低精度需求已从 FP8 矩阵计算进一步扩展至 FP4 计算与存储。与此同时，模型各部分对数值误差的容忍程度不同，输出层、注意力等敏感计算仍需保留 BF16 等较高精度的矩阵运算~\cite{deepseekv3}。因此，实际负载对张量计算硬件提出了多种精度并存的需求。但是，从低精度浮点本身特性来看，低位宽能够减少数据存储与搬运开销，并为在有限硬件资源下提高计算并行度提供空间，成为未来AI模型计算精度选择的一个重要方向。而在面积与功耗预算约束下，如何兼顾较高精度计算能力与低精度吞吐需求，进而确定 16、8、4 位浮点张量计算的峰值算力配比，也成为当前乃至未来AI芯片的重要架构设计问题。

当前产业界的配比选择已呈现多样化趋势。以同一芯片的 16 位峰值算力归一化，按各精度稠密矩阵计算峰值比较，NVIDIA B200、AMD MI355X 与华为 Ascend 950PR 采用经典的 $1:2:4$ 配比~\cite{nvidiahgx,amdmi355x,huawei950}；NVIDIA B300 将配比调整为 $1:2:6$~\cite{nvidiahgx}，AMD MI455X 约为 $1:4:8$~\cite{amdmi455x}，华为 Ascend 960PR 的公开研发路线目标则为 $1:2:8$~\cite{huawei960roadmap}。从硬件规划的角度看，低精度算力增配具有面向未来负载的前瞻性。芯片从架构设计、验证到流片部署需要较长周期，而模型结构与量化方法迭代更快；设计阶段观察到的精度需求，未必能代表芯片部署及使用期间的负载。因此，配比选择需要预估未来低精度计算需求的增长，在芯片设计阶段提前配置相应资源。

然而，公开规格展示了配比选择的结果，却不足以说明这些配比如何实现，以及新增的面积、功耗与带宽需求能够换来多少可用算力。即使配比相同，不同的归约长度、输出并行度与后端共享范围，也会改变资源开销及算力发挥所需的数据供给条件。现有多精度计算与设计空间探索研究已考察通路复用、数值格式及硬件成本之间的关系~\cite{transdot2026,paretomac2025,zhao2026mxlike}，但面向算力配比及其实现方式，仍缺少可复现、可对比的系统评估。面向尚未确定的未来负载，需要将配比、实现方式与资源约束纳入统一框架，考察不同精度需求下的成本与算力收益，为前瞻性配比选择提供依据。

为此，本文提出 DP-Puzzle，一种面向多精度浮点张量计算的分层可组合点积架构。在点积单元（dot-product unit）内部，通过跨精度资源复用与低精度专用计算模块的组合，沿归约维度 $K$ 增加目标精度下每周期处理的乘积项数；在阵列层面，通过组合具有不同精度支持与计算能力的点积单元，沿输出维度 $N$ 调整各精度下的输出并行度。两类组合分别在单元内部与阵列层面改变算力配比。在给定 Tensor Core 或阵列面积预算下，本文结合面积、功耗与带宽模型，比较两类实现，并分析不同精度负载及数据供给条件下的算力收益与帕累托权衡，为未来芯片的配比规划提供可比较的硬件实现与成本依据。

本文的主要贡献如下：

\begin{enumerate}
  \item \textbf{支持多精度浮点精度的点积单元设计。} 结合新型浮点精度格式定义，通过合理设计分组，针对乘法器，加法器与移位器资源复用，支持BF16、FP16、MXFP8/6/4、NVFP4与HiF4等多精度浮点计算。
  \item \textbf{提出两种改变传统配比的实现方式，并进行硬件成本与收益分析。} 在匹配各精度绝对计算能力的条件下，对比单元内 $K$ 向组合与阵列级 $N$ 向组合的面积、功耗与带宽成本。在所评估配置中，单元内 $K$ 向扩展具有面积优势，能耗收益则与目标精度和阵列规模有关；两类组合的带宽需求增长随阵列输出维规模而变化。
  \item \textbf{建立负载驱动的配比选择方法。} 针对四种代表配比，建立固定整核面积、整数部署和混合精度时间／能耗模型，以精确选择边界连接硬件成本与真实模型结构导出的 prefill 工作量，揭示配比收益随上下文和执行精度变化的规律。
\end{enumerate}

本文其余部分安排如下：第~\ref{sec:background} 节介绍背景与相关工作；第~\ref{sec:architecture} 节说明 DP-Puzzle 架构及配比实现方法；第~\ref{sec:evaluation} 节评估组合成本与配比收益；第~\ref{sec:conclusion} 节总结全文。
